\chapter{Proposed Solution: Workflow-First Deployment via Function Registry Resolution}
\label{ch:proposed_solution}

This chapter introduces the core idea of the proposed OmniFlow-QuickFaaS integration. The objective is to reduce the manual work required to deploy workflows that depend on serverless functions, while keeping the workflow definition readable and stable.

In the reference OmniFlow usage, a \texttt{Call} step is expressed as an HTTP request: the developer supplies concrete invocation details such as \texttt{method}, \texttt{host} and \texttt{path}. As a result, whenever the workflow calls a serverless function, the developer must already know its endpoint, which exists only once the function is deployed and differs for each account, project and provider. The proposed solution removes this burden by letting the developer declare, directly on the call step, an endpoint that is resolved automatically: a dedicated \texttt{internalFunction(name, deploymentDescriptorPath)} construct replaces the plain \texttt{host}/\texttt{path} fields used to provide a literal endpoint, the two being mutually exclusive. This enables a single, workflow-first deployment that resolves (or, if needed, provisions) the required functions and then deploys the workflow to the chosen cloud provider.

The chapter is organised in four sections. Section~\ref{sec:components-architecture} presents the
components that the integration adds to the toolchain, how they cooperate at deployment time, and
the three-level resolution cascade (\S\ref{sec:resolution-cascade}) that decides, for each call,
whether a function is looked up, discovered or deployed. Section~\ref{sec:design-rationale}
justifies the two main design decisions behind that arrangement: why OmniFlow drives QuickFaaS
rather than the reverse, and why the two tools are combined in a single deployment process instead
of running as independent communicating services. Section~\ref{sec:dsl-extensions} describes the
extension made to the OmniFlow DSL, the \texttt{internalFunction} construct on the call step,
together with the concepts it relies on: function ownership, the \texttt{functionRef} identifier,
the Function Registry and the QuickFaaS deployment descriptor. Section~\ref{sec:developer-workflow}
closes the chapter from the developer's point of view, with the artifacts to prepare, a complete
example workflow, and a walkthrough (\S\ref{sec:walkthrough}) of what the integrated deployment
does on a first and on a second run, including the failure modes it reports.

\section{Components Architecture}
\label{sec:components-architecture}

The proposed integration keeps the original OmniFlow workflow model but extends the toolchain with a small set of additional components that cooperate at deployment time.

Figures~\ref{fig:architecture-before} and~\ref{fig:architecture-after} contrast the developer workflow before and after the integration proposed in this dissertation.

In the prior state (Figure~\ref{fig:architecture-before}), OmniFlow and QuickFaaS are used independently. The developer first deploys each function with QuickFaaS in a separate manual step, retrieves the resulting invocation endpoint, and copies it by hand into the workflow definition before deploying the workflow: the manual synchronisation described in \S\ref{sec:int_cha_int_uni_multi}, with the ordering it imposes and the errors it invites. Moreover, QuickFaaS targeted only GCP and Azure, so along the AWS path OmniFlow had no function-deployment capability at all. The two lifecycle concerns, function deployment and workflow deployment, are therefore completely disconnected.

\begin{figure}[htbp]
  \centering
  \includegraphics[width=0.85\textwidth]{images/architecture/architecture-before.png}
  \caption{OmniFlow and QuickFaaS operate as two disjoint tools.}
  \label{fig:architecture-before}
\end{figure}

In the state introduced by this work (Figure~\ref{fig:architecture-after}), a single workflow definition drives both concerns. Each provider deployer delegates call resolution to a provider-specific resolver — \texttt{Workflow\allowbreak Internal\allowbreak Function\allowbreak Resolver} for GCP and \texttt{Aws\allowbreak Internal\allowbreak Function\allowbreak Resolver} for AWS — and each is backed by its own registry file (\texttt{function-registry.gcp.json} and \texttt{function-registry.aws.json} by default) that maps logical identifiers to concrete invocation metadata. QuickFaaS is extended with an \texttt{AwsProvider}, so a function can be deployed on demand to AWS Lambda through the same subprocess mechanism already used for GCP. Once resolved, the rendered Step Functions state machine invokes the Lambda natively by ARN. The two halves of the contribution are visible in a single view: the AWS provider inside QuickFaaS and the resolver–registry layer that unifies function and workflow deployment.

\begin{figure}[htbp]
  \centering
  \includegraphics[width=\textwidth]{images/architecture/architecture-after.png}
  \caption{Unified deployment driven from a single workflow definition; components marked as new/extended are the contribution of this work.}
  \label{fig:architecture-after}
\end{figure}

The end-to-end flow proceeds as follows:
\begin{enumerate}
  \item the developer defines the workflow and one QuickFaaS descriptor per function to be
  resolved this way;
  \item OmniFlow parses the workflow and, for each such call,
  \item resolves the endpoint from the registry, discovers it directly on the provider, or
  deploys it through QuickFaaS;
  \item metadata obtained by discovery or deployment, or refreshed after drift, is written to
  the registry; and
  \item the resolved workflow is rendered and deployed to the provider's workflow service.
\end{enumerate}
Calls that already carry a literal endpoint bypass the registry entirely.

\subsection{The Resolution Cascade}
\label{sec:resolution-cascade}

Step~(3) of the end-to-end flow listed in Section~\ref{sec:components-architecture} is the heart
of the integration and is best understood as a three-level cascade,
applied independently to each call resolved this way and stopping at the first level that succeeds. The
cascade is designed so that the cheapest and least intrusive action is always preferred, and an
actual deployment is a last resort rather than a default.

\begin{enumerate}
  \item \textbf{Level~1 (Registry).} The registry is consulted for the \texttt{functionRef}
  (exact match first, then a suffix match on \texttt{"region/functionRef"}). A hit is the common
  case for a function that has already been deployed once. The stored endpoint is then validated
  against the live resource (the Cloud Run service on GCP, the Lambda function on AWS) so that
  a stale binding cannot be used silently: a changed endpoint refreshes the entry, and a deleted
  function is looked for in the other regions, its entry replaced if it is found there and removed
  otherwise, in which case the deployment aborts. Validation contacts the provider, but the registry keeps it
  cheap: an entry records where the function lives (the region inside the Lambda ARN, the full
  resource name of the Cloud Run service), so validation is one lookup in a known place rather than
  the search across every region that Level~2 performs. On GCP, bindings to the first-generation
  Cloud Functions that QuickFaaS deploys skip this check (\S\ref{subsec:store-internals}).
  \item \textbf{Level~2 (Provider discovery).} On a registry miss, the target provider is queried
  directly (a Cloud Run lookup across regions on GCP; on AWS, a \texttt{GetFunction} lookup across
  every region enabled for the account, or a single call when the caller supplies the
  fully-qualified \texttt{"region/functionRef"} form). If the
  function already exists (for instance, because it was deployed outside the integration), its
  metadata is recorded in the registry and used, again \emph{without} triggering a deployment.
  \item \textbf{Level~3 (QuickFaaS deployment).} Only when the function is absent from both the
  registry and the provider, and the call supplied a \texttt{deploymentDescriptorPath}, is QuickFaaS
  invoked to deploy it; the returned metadata is then registered and used. If no descriptor was
  provided, deployment fails fast with a diagnostic naming the unresolved function.
\end{enumerate}

This ordering gives the integration three desirable properties. It is \emph{idempotent} at the
level that matters: re-deploying the same workflow re-resolves from the registry (or the
provider) and does not redeploy an unchanged function. It is \emph{safe by default}: validation
and discovery never mutate the provider, and a deployment happens only on an explicit descriptor
and only after the function's absence is confirmed. A deployment does change more than the
function itself: it also adjusts permissions: on GCP, when an invoker service account is
configured, by granting it permission to invoke the function; on AWS, by creating or reusing an
execution role, repairing the trust policy of a role the developer supplies, and granting Step
Functions permission to invoke the function (\S\ref{subsec:omniflow-aws-deployer}). And it is
\emph{observable}: each level emits an explicit diagnostic, so the developer can see whether a
call was served from the registry, discovered on the provider, or freshly deployed. The concrete
class-level realisation of this cascade, including drift handling and ambiguity detection, is
described in Chapter~\ref{cha:impl}.

\paragraph{Static binding.}
The cascade runs only when the workflow is deployed. The endpoint it resolves is written into the
rendered workflow (as the \texttt{url} of an HTTP call in a GCP Workflows definition, and as the
\texttt{FunctionName} ARN of a \texttt{lambda:invoke} task in a Step Functions state machine),
and the deployed workflow never consults the registry at run time. The binding is therefore
\emph{static}: if a function's endpoint changes after the workflow has been deployed, the deployed
workflow keeps calling the old endpoint until it is deployed again. On that next deployment,
Level~1 validation detects the change and binds the call to the new endpoint, except for GCP
first-generation Cloud Functions, whose bindings are not validated. In practice an endpoint changes
less often than a function's code: updating a Lambda's code in place keeps its ARN, and
redeploying a Cloud Function or Cloud Run service under the same name and region keeps its URL. It
changes when the function is recreated under another name, region, account or project. Resolving
endpoints at run time instead would require every workflow execution to look the function up,
which is outside the scope of this work.

\paragraph{Provider and region scope.}
Cloud providers organise their infrastructure into geographic \emph{regions} (e.g., \texttt{eu-west-1}
on AWS or \texttt{europe-west1} on GCP), and every resource (a Lambda function, a Cloud Run
service, a storage bucket, or the deployed workflow itself) lives in a specific region, which is
part of its identity (on AWS the region even appears inside the function ARN). This work is scoped
to a \emph{single provider per deployment}: a workflow deployed to GCP resolves only GCP functions,
and one deployed to AWS resolves only AWS Lambdas. Within a provider, the workflow and its registry-resolved
functions are \emph{not} strictly required to share a region, and on neither provider is the
developer required to declare one: Level~2 searches every region (on AWS, every region the EC2
\texttt{DescribeRegions} API reports as enabled for the account), and the fully-qualified
\texttt{"region/functionRef"} form serves to disambiguate a name found in more than one region.
On both providers the search starts with the region the workflow is being deployed to, which
reflects the expected deployment, in which the
workflow and its functions share a region: the resolver does not require it, but same-region
deployment minimises invocation latency and cross-region data-transfer cost, and simplifies
data-residency compliance. One co-location constraint is, however, mandatory
on AWS: a Lambda must reside in the same region as the S3 bucket that holds its deployment package,
so the AWS deployer derives the effective region from that bucket
(Chapter~\ref{cha:impl}).

\section{Design Rationale: Embedding Function Deployment in OmniFlow}
\label{sec:design-rationale}

A central design decision of this integration is that OmniFlow \emph{drives} QuickFaaS as part of a
single workflow deployment, rather than the two tools being deployed as independent, long-running
services that communicate over a network API. This section justifies that choice and states its
trade-offs explicitly.

\paragraph{Why OmniFlow orchestrates QuickFaaS.}
Only OmniFlow has the information needed to drive a \emph{workflow-first} deployment: it owns the
deployment entry point and the workflow tree, and therefore knows which functions the
workflow references and in what order they must be resolved. QuickFaaS, by contrast, operates at the
granularity of a single function described by a deployment descriptor and has no notion of a
workflow. The orchestration must therefore live on the OmniFlow side, with QuickFaaS acting as the
function-deployment mechanism it invokes on demand. Crucially, this is a \emph{deployment-time},
client-side, one-shot action performed on the developer's machine, not a recurring runtime
interaction between two systems.

To keep the two tools decoupled despite this dependency, function deployment is reached through the
\texttt{Internal\allowbreak Function\allowbreak Deployer} strategy interface, whose provider implementations
(\texttt{QuickFaasDeployer} for GCP, \texttt{AwsLambdaDeployer} for AWS) invoke QuickFaaS as an
external process. QuickFaaS is thus reused as-is, remaining a separately owned and separately built
tool (its own repository, Gradle build and Kotlin~1.6.20 toolchain), instead of having its source
merged into OmniFlow's Maven build.

\paragraph{Why not two independent communicating services.}
An alternative design would run OmniFlow and QuickFaaS as two standing services that exchange
messages over a network. This was rejected for several reasons. First, neither tool is a
long-running service: OmniFlow is a client-side library invoked during a build, and QuickFaaS is a
self-contained desktop application that consumes a local deployment descriptor and exposes no
network API. Turning either into a service would require re-architecting it and then hosting,
securing and keeping it available. Second, doing so would reintroduce precisely the
``additional runtime layer that must be operated'' which this thesis deliberately avoids
(Section~\ref{ch:related-works}): the contribution is aligned with provider-native execution
surfaces and adds no new infrastructure of its own. Third, deployment uses the developer's own
cloud credentials, resolved locally (e.g., application-default credentials or environment
variables); a service-to-service design would have to transport or delegate those credentials
across a trust boundary, a security liability with no corresponding benefit for a single-developer,
single-machine build. Finally, for a task that is inherently sequential, local and ephemeral, a
direct invocation has far fewer failure modes than a distributed protocol (no network partitions,
partial failures, or inter-service API to version), and the coupling is in fact looser: OmniFlow
depends on the QuickFaaS build artifact, not on a running service.

\paragraph{Trade-off of the current mechanism.}
In the present prototype the invocation is realised as a subprocess: OmniFlow places the deployment
descriptor (and the function source) next to the QuickFaaS JAR and launches it with
\texttt{java\ -jar}, detecting success from the process exit code and a marker printed on its
output. This is pragmatic (it reuses QuickFaaS's existing headless entry point without modifying
it), but it is also the most brittle part of the integration, since it couples the two build
systems at run time through a JAR path and parses the child process's output. Importantly, the
cleaner alternative to the subprocess is \emph{in-process library invocation} (linking QuickFaaS as
a dependency), which preserves the OmniFlow-drives-QuickFaaS model while removing the process
boundary; this is discussed as future work (Chapter~\ref{cha:conclusions}) and is a different axis
from the two-service design rejected above.

\section{DSL Extensions}
\label{sec:dsl-extensions}

The OmniFlow background chapter described how a workflow is modelled as a tree of \texttt{Step} nodes, each with a single \texttt{context} that determines its behaviour (Call, Assign, Switch, Iteration, Parallel). Within this hierarchy, \emph{execution steps} correspond to nodes that actually perform work (such as HTTP calls or variable assignments). This section focuses on one particular execution step, the \texttt{call} step, and on how it is extended in this dissertation.

Figure~\ref{fig:call-step-model} shows how the \texttt{call} step is extended to express calls whose endpoint is resolved automatically. In the original design, a \texttt{CallContext} carries the literal HTTP fields \texttt{host} and \texttt{path}. This work adds an optional association to an \texttt{InternalFunction}, which holds a logical \texttt{name} (the \texttt{functionRef}) and an optional \texttt{deploymentDescriptorPath}. The two construction paths are \emph{mutually exclusive}: a call with a literal endpoint provides \texttt{host}/\texttt{path} and no \texttt{internalFunction}; a call resolved through the mechanism described in this chapter provides an \texttt{internalFunction} and \emph{must not} set \texttt{host} or \texttt{path}, since those are injected automatically at deployment time by the resolution mechanism. Keeping this form as an alternative on the existing step, rather than introducing a new step type, means developers continue to describe calls with the familiar HTTP vocabulary.

\begin{figure}[htbp]
  \centering
  \includegraphics[width=0.8\textwidth]{images/dsl-extension/call-step-model.png}
  \caption[Extended model of the \texttt{call} execution step.]{Extended model of the
  \texttt{call} execution step, with the mutually exclusive \texttt{internalFunction} association
  added in this work.}
  \label{fig:call-step-model}
\end{figure}

With this extension, the call step can be used in two distinct ways:

\begin{itemize}
  \item With a \textbf{literal endpoint}, the call step is used exactly as in the original OmniFlow design, and the integration triggers no function deployment.
  \item With an \textbf{automatically resolved endpoint}, the endpoint is resolved at deployment time, and the function is first deployed from its descriptor if it does not exist.
\end{itemize}

The concepts this second form relies on, and through which it connects to the QuickFaaS-based
deployment and registry mechanisms described in the following sections, can be summarised as
follows:

\paragraph{Function Ownership Determines Resolution.}
Which of the two forms a call uses depends on who owns the target. A function whose source code the workflow developer owns, and is allowed to deploy and update, is declared through \texttt{internalFunction}. A function or API the developer does not own (for example, a third-party service) is called through a literal \texttt{host}/\texttt{path}, and is treated as an unmanaged HTTP endpoint: it is never deployed or updated through QuickFaaS and is not recorded in the registry.

\paragraph{Function Reference (\texttt{functionRef}).}
A function reference (\texttt{functionRef}) is a stable logical identifier used to name a function whose endpoint is resolved through the registry (e.g., \texttt{"fraud-check"}). This identifier is used consistently across the workflow definition, the registry metadata and the deployment descriptors. It is not an endpoint, it is a name that allows OmniFlow to bind a call to the concrete endpoint during deployment.

\paragraph{Function Registry.}
The function registry is a file, one per provider, that contains invocation metadata only for functions resolved this way. Conceptually, it maps each such \texttt{functionRef} to the concrete values required to invoke it.

Listing~\ref{lst:registry-structure} shows the structure of the file, with one example per
provider. The document has two top-level fields. \texttt{updatedAt} records, as an ISO-8601
timestamp, when the file was last written, so that a developer inspecting it can tell how recent
the bindings are. \texttt{functions} is an object used as a map: each key is a
\texttt{functionRef} and the corresponding value is the invocation metadata bound to it. The key is
the bare reference (\texttt{"fraud-check"} in the GCP example) while it is unambiguous, and is
promoted to the qualified \texttt{"region/functionRef"} form when the same name is found in more
than one region, which is what allows an ambiguous reference to be reported instead of being
resolved arbitrarily. Each value carries exactly two fields: \texttt{serviceName}, the name the
function has on the provider, and \texttt{url}, the single value the rendered workflow needs in
order to invoke it.

The two example files differ only in that last field, and this is where the provider shows through.
On GCP, \texttt{url} is the HTTPS endpoint of the function
(\texttt{https://\allowbreak europe-west1-\allowbreak proj.\allowbreak cloudfunctions.\allowbreak net/\allowbreak fraud-check}), which the resolver
splits into the \texttt{host} and \texttt{path} of an ordinary HTTP call. On AWS, \texttt{url} is
the function's ARN
(\texttt{arn:aws:\allowbreak lambda:\allowbreak eu-west-1:\allowbreak 123456789012:\allowbreak function:\allowbreak hello-lambda-fn}),
which the rendered state machine uses as the resource of a native Lambda invocation rather than as
an HTTP address. Nothing else is
stored: no credentials and no deployment parameters. Both files are written and refreshed by the integration, never by hand.

\begin{lstlisting}[caption={Example registry files, one per provider (registry-managed functions only).},label={lst:registry-structure},basicstyle=\ttfamily\footnotesize,breaklines=false]
function-registry.gcp.json:
{
  "updatedAt": "2026-02-12T18:40:00Z",
  "functions": {
    "fraud-check": {
      "serviceName": "fraud-check",
      "url": "https://europe-west1-proj.cloudfunctions.net/fraud-check"
    }
  }
}

function-registry.aws.json:
{
  "updatedAt": "2026-02-12T18:40:00Z",
  "functions": {
    "hello-lambda-fn": {
      "serviceName": "hello-lambda-fn",
      "url": "arn:aws:lambda:eu-west-1:123456789012:function:hello-lambda-fn"
    }
  }
}
\end{lstlisting}

Listing~\ref{lst:call-step-internal} shows a complete call step whose endpoint is resolved through
the registry.

\begin{lstlisting}[language=Kotlin,caption={Call step declared via \texttt{internalFunction}.},label={lst:call-step-internal}]
step {
  name("call-auto-deployed-function")
  description("Calls a function that QuickFaaS will deploy if not already running")
  context(
    call {
      method(GET)
      internalFunction(
        "quickfaas-test-fn",
        deploymentDescriptorPath = "./functions/quickfaas-test-fn/func-deployment.json"
      )
      result("result")
    }
  )
}
\end{lstlisting}

The step is written with the vocabulary of any other OmniFlow step: \texttt{name} identifies it
inside the workflow, \texttt{description} documents it, and \texttt{context(call\{\dots\})} makes it
a call step. Inside the call, \texttt{method(GET)} still declares the HTTP verb, and
\texttt{result("result")} names the workflow variable that receives the response. The only
difference from a conventional call is what is \emph{not} there: no \texttt{host} and no
\texttt{path}. In their place, \texttt{internalFunction} supplies the \texttt{functionRef}
\texttt{"quickfaas-test-fn"}, the key under which the endpoint is stored in (or added to) the
registry, and \texttt{deploymentDescriptorPath}, a path to the QuickFaaS descriptor that is read
only if the cascade reaches Level~3 and the function has to be deployed. Writing the step this way
is a statement about ownership rather than about addressing: it says which function to call, and
leaves it to the deployment to determine where that function is.

\paragraph{QuickFaaS Deployment Descriptor and Code File}
To deploy a function managed this way, QuickFaaS requires two artifacts: (i) a JSON
deployment descriptor and (ii) the function code file. Listing~\ref{lst:quickfaas_json_config}
in Section~\ref{sec:background_quickfaas} is an example of such a descriptor, and its fields fall
into three groups.

The first group says \emph{where} the function is to be deployed: \texttt{cloudProvider} selects
the target platform (\texttt{"gcp"} or, through the provider added in this work, \texttt{"aws"}),
\texttt{project} identifies the account under which the function is created (a GCP project
identifier, or the AWS account number), and \texttt{accessToken} carries the credential QuickFaaS
authenticates with on GCP; on AWS the field is left empty, because the provider added in this
work takes its credentials from the environment instead. The second group, nested under \texttt{function}, says \emph{what} is created:
\texttt{name} is the name the function will have on the provider, \texttt{location} its region,
\texttt{runtime} the language runtime (for example \texttt{"java17"}), \texttt{trigger} the kind of
event that invokes it (\texttt{"http"} for the functions a workflow calls), and \texttt{bucket} the
object-storage bucket through which the deployment package is staged. The third group is the code
itself: \texttt{functionFile} is a filesystem path to the cloud-agnostic implementation described
in Section~\ref{sec:background_quickfaas}, which QuickFaaS packages and uploads.

Two of these fields matter to the integration beyond deployment. The \texttt{function.name} must
equal the \texttt{functionRef} used on the call step, because the resolver looks the deployed
function up under the reference it was given; the two names are not checked against each other, and
the consequence of letting them diverge is described at the end of
\S\ref{sec:walkthrough}. And \texttt{function.location} fixes the region the
resulting registry entry will point at, which is what a later resolution validates against. The
descriptor is otherwise the same artifact QuickFaaS has always taken, read now by OmniFlow instead
of by the developer running QuickFaaS by hand; the AWS variant, which adds an \texttt{iamRoleArn}
field for the Lambda execution role, is shown in
Listing~\ref{lst:aws-descriptor}.

\section{Developer Procedure for Deploying a Workflow}
\label{sec:developer-workflow}

To deploy a workflow that introduces one new function to be resolved through the registry and calls one function invoked through a literal endpoint, the developer prepares a QuickFaaS deployment descriptor file (JSON) for the former and the corresponding code file referenced by \texttt{functionFile}. No QuickFaaS descriptor is created for the latter, since the developer is not its owner and the integrated tool must not deploy or update it.

Because such calls are declared with \texttt{internalFunction("fraud-check", ...)} rather than a concrete endpoint, the developer only needs to choose a stable \texttt{functionRef}. Read access to the Function Registry is useful to discover which functions are already registered and which \texttt{functionRef} identifiers are already in use, but it is not required for a brand-new function: after the first integrated deployment, OmniFlow records the binding so the same \texttt{functionRef} can be reused reliably in future workflows.

The workflow is written as shown in Listing~\ref{lst:wf-internal-part} and Listing~\ref{lst:wf-external-part}. Listing~\ref{lst:wf-internal-part} shows the first step, which calls a function resolved through the registry: instead of a concrete endpoint it uses \texttt{internalFunction(\allowbreak "fraud-check",\allowbreak "./deploy/fraud-check.json")}, whose second argument points to the QuickFaaS descriptor. Listing~\ref{lst:wf-external-part} completes the workflow with a second step that calls an API using literal \texttt{host}/\texttt{path} values and no \texttt{internalFunction}. What the integrated deployment does with each of the two calls is traced in \S\ref{sec:walkthrough}.

\begin{lstlisting}[language=Kotlin,caption={FraudPipeline workflow (part 1/2): call resolved through the registry.},label={lst:wf-internal-part},basicstyle=\ttfamily\small,breaklines=false]
val wf = workflow {
  name("FraudPipeline")
  input("args")
  steps(
    step {
      name("fraud-check")
      context(call {
        method(GET)
        internalFunction("fraud-check", "./deploy/fraud-check.json")
        result("fraudResult")
      })
    },
    ...
\end{lstlisting}

\begin{lstlisting}[language=Kotlin,caption={FraudPipeline workflow (part 2/2): call with a literal endpoint.},label={lst:wf-external-part},basicstyle=\ttfamily\small,breaklines=false]
...
// FraudPipeline workflow (continuation)
    step {
      name("external-risk-api")
      context(call {
        method(POST)
        host("api.partner.example")
        path("/v1/risk-score")
        result("riskResult")
      })
    }
  )
  result("riskResult")
}
\end{lstlisting}

\subsection{Walkthrough of an Integrated Deployment}
\label{sec:walkthrough}

Deploying the \texttt{FraudPipeline} of Listings~\ref{lst:wf-internal-part}
and~\ref{lst:wf-external-part} makes the cascade concrete. Suppose \texttt{fraud-check} has never
been deployed and the registry does not yet exist. OmniFlow first bootstraps the registry (on GCP by listing the project's HTTP-invocable
services, on AWS by listing the Lambda functions across every region enabled for the account) and then walks the workflow tree. The \texttt{external-risk-api} step carries literal \texttt{host}/\texttt{path} and
is skipped by the resolver. The \texttt{fraud-check} step carries an \texttt{internalFunction}, so
the cascade runs: Level~1 misses (empty registry), Level~2 misses (the function does not exist on
the provider), and Level~3 fires because the call supplied
\texttt{"./deploy/fraud-check.json"}. QuickFaaS deploys the function, the returned endpoint (an
HTTPS URL on GCP, an ARN on AWS) is written to the registry under the key \texttt{fraud-check}, and
the resolved value is injected into the step. With every such call now bound, OmniFlow renders
the provider artifact and deploys the workflow.

A second deployment of the same workflow behaves very differently and illustrates the value of the
indirection: \texttt{fraud-check} is now a Level~1 registry hit, no descriptor is consulted, and no
function is redeployed. The workflow definition did not change, and it would not have to change if,
in the meantime, \texttt{fraud-check} had been recreated at a new endpoint (for example, in another
region): as the static binding of \S\ref{sec:resolution-cascade} describes, this deployment would
bind the call to the new endpoint, while the workflow deployed earlier keeps calling the old one.

\paragraph{Error and edge cases.}
The same procedure defines a small number of failure modes, all fail-fast and all reported before
any provider artifact is rendered:
\begin{itemize}
  \item \textbf{Unresolvable function.} Absent from the registry and the provider, with no
  descriptor supplied: deployment aborts, naming the function and the two corrective actions
  (deploy it first, or add a \texttt{deploymentDescriptorPath}).
  \item \textbf{Ambiguous reference.} A bare \texttt{functionRef} that matches entries in more
  than one region (in the registry or on the provider): deployment aborts and asks the developer to
  use the fully-qualified \texttt{"region/functionRef"} form.
  \item \textbf{Mutually exclusive forms.} A call that sets both \texttt{internalFunction} and
  \texttt{host}/\texttt{path}: rejected at DSL build time and again at resolution time.
  \item \textbf{Stale binding.} A registry entry whose function no longer exists (a deleted
  Lambda function on AWS, a deleted Cloud Run service on GCP) and cannot be rediscovered in another
  region: the entry is removed and deployment aborts, even when the call supplies a descriptor,
  rather than emit a workflow bound to a dead endpoint. A stale entry for a first-generation Cloud
  Function on GCP goes undetected, since those entries are not validated.
\end{itemize}

One misconfiguration is not caught before deployment: a \texttt{functionRef} that differs from the
\texttt{function.name} in its descriptor. The resolver passes the \texttt{functionRef} to the
deployer as the function's name, while QuickFaaS deploys the function under the descriptor's name,
and nothing checks that the two agree. When they differ, OmniFlow waits for a function named after
the \texttt{functionRef}, and the deployment fails when that wait times out; the function that was
deployed is not recorded in the registry.











